\chapter[\emj{🔤}~Tries (Árboles de Prefijos)]{\emj{🔤}~Tries (Árboles de Prefijos)}

\begin{dsaquees}

Un diccionario físico 📖 donde navegas letra por letra. Para buscar
la palabra "GATO":
1. Vas a la sección G.
2. Dentro de G, buscas la A.
3. Dentro de GA, buscas la T.
4. Dentro de GAT, buscas la O.
¡Encontrado!

Cada nivel del árbol representa una letra. Si compartes el mismo
inicio (prefijo), compartes los mismos nodos. Por ejemplo, "GATO" y 
"GOTA" comparten la 'G', pero luego se bifurcan.

\end{dsaquees}

\begin{dsaotraforma}

Un Trie (se pronuncia como "try") es un árbol especializado para
guardar strings. Cada nodo representa un caracter. 

Su gran ventaja es la eficiencia para búsquedas de prefijos y 
autocompletado. A diferencia de un Hash Map, el tiempo de búsqueda
en un Trie depende únicamente de la longitud de la palabra (L), no
del número de palabras guardadas (N).

\end{dsaotraforma}

\begin{dsadiagrama}
\begin{verbatim}
Insertando: "CAT", "CAR", "DO", "DOG"

      [root]
      /    \
    [C]    [D]
     |      |
    [A]    [O]*  <-- "DO" es palabra (isEnd=true)
    / \     |
  [T]*[R]* [G]*  <-- "CAT", "CAR", "DOG" son palabras

* : Indica el final de una palabra (isEnd = true).
\end{verbatim}
\end{dsadiagrama}

\begin{lstlisting}
1. insert(word): Empiezas en el root. Para cada letra, si no existe el
   hijo, lo creas. Al final de la palabra, marcas `isEnd = true`.
2. search(word): Sigues el camino de las letras. Si llegas al final
   y el nodo tiene `isEnd = true`, la palabra existe.
3. startsWith(prefix): Igual que search, pero no importa si `isEnd` 
   es true; con que el camino exista, el prefijo es válido. O(L).

PSEUDOCÓDIGO
──────────────────────────────────────────────────────────────

// Insertar palabra en el Trie
función insertar(palabra):
    actual = root
    PARA cada letra EN palabra:
        SI actual.hijos[letra] NO existe:
            actual.hijos[letra] = nuevo Nodo()
        actual = actual.hijos[letra]
    actual.esFinDePalabra = VERDADERO

CÓDIGO C++
──────────────────────────────────────────────────────────────

#include <iostream>
#include <unordered_map>
#include <string>
using namespace std;

struct TrieNode {
    unordered_map<char, TrieNode*> children;
    bool isEndOfWord = false;
};

class Trie {
private:
    TrieNode* root;

public:
    Trie() { root = new TrieNode(); }

    // Insertar -> O(L) donde L es la longitud de la palabra
    void insert(string word) {
        TrieNode* curr = root;
        for (char c : word) {
            if (curr->children.find(c) == curr->children.end()) {
                curr->children[c] = new TrieNode();
            }
            curr = curr->children[c];
        }
        curr->isEndOfWord = true;
    }

    // Buscar palabra completa -> O(L)
    bool search(string word) {
        TrieNode* curr = root;
        for (char c : word) {
            if (curr->children.find(c) == curr->children.end()) return false;
            curr = curr->children[c];
        }
        return curr->isEndOfWord;
    }

    // Buscar prefijo -> O(L)
    bool startsWith(string prefix) {
        TrieNode* curr = root;
        for (char c : prefix) {
            if (curr->children.find(c) == curr->children.end()) return false;
            curr = curr->children[c];
        }
        return true;
    }
};

int main() {
    Trie t;
    t.insert("gato");
    t.insert("goma");

    cout << "¿Existe gato? " << t.search("gato") << endl;    // 1 (true)
    cout << "¿Existe prefijo go? " << t.startsWith("go") << endl; // 1 (true)
    cout << "¿Existe gatito? " << t.search("gatito") << endl; // 0 (false)

    return 0;
}

COMPLEJIDAD (Big-O)
──────────────────────────────────────────────────────────────

  Operación        │ Tiempo    │ Espacio
  ─────────────────┼───────────┼──────────────────────────────
  Insert           │ O(L)      │ O(ALPHABET_SIZE * L * N)
  Search           │ O(L)      │ O(1)
  StartsWith       │ O(L)      │ O(1)

* L = Longitud de la palabra.
* N = Número de palabras.
* Nota: El espacio puede ser grande si las palabras no comparten prefijos.
\end{lstlisting}

\begin{dsaentrevistas}

✅ Cuándo usar: Si el problema menciona "autocompletar", "diccionario",
   o "buscar todas las palabras que empiecen con...", usa un Trie.

💡 Trie + DFS: Una combinación clásica para problemas como 
   "Word Search II" (encontrar palabras en una cuadrícula de letras).

⚠️ Hash Map vs Trie: 
   - Hash Map: Busca la palabra completa en O(1) promedio.
   - Trie: Puede buscar prefijos y es más rápido si hay muchas colisiones
     en el hash map.

🔑 Memoria: Si te preocupa el espacio, menciona que se puede usar un
   "Compressed Trie" (Radix Tree) donde los nodos con un solo hijo 
   se fusionan.

\end{dsaentrevistas}

\begin{dsaresumen}

• Árbol de letras, nivel por nivel.
• Ideal para búsquedas de prefijos y autocompletado.
• Tiempo de búsqueda dependiente solo del largo de la palabra.
• `startsWith` es su operación estrella.
• Puede usar mucha memoria, pero ahorra espacio si hay muchos 
  prefijos compartidos.
════════════════════════════════════════════════════════════════

\end{dsaresumen}


